\chapter{The Game Loop}
\label{ch:game-loop}

Chapter~\ref{ch:first-game} showed \texttt{Update} / \texttt{Draw} as two overridden methods
without asking what drives them. This chapter opens up \cnaclass{Game} itself --- the class
every CNA game inherits from --- and the small constellation of supporting types
(\cnaclass{GameTime}, \cnaclass{GameComponent}, \cnaclass{GameWindow},
\cnaclass{GameServiceContainer}) that make the loop work.

\section{Game: base class, not framework-owned entry point}

\cnaclass{Game} inherits \texttt{System::Object} and \texttt{System::IDisposable}, and its
own Doxygen brief describes it precisely: ``base class that provides the XNA-style game
loop, services, window, content and components.'' Four events --- \texttt{Activated},
\texttt{Deactivated}, \texttt{Disposed}, \texttt{Exiting} --- are all typed
\texttt{System::EventHandler<System::EventArgs>}. This is worth pausing on for
\texttt{Exiting} specifically: CNA does declare a dedicated \cnaclass{ExitingEventArgs} class
(matching real XNA's public API surface), but \texttt{Game::Exiting} itself is typed against
plain \texttt{EventArgs}, not \texttt{EventHandler<ExitingEventArgs>} --- and grepping the
tree confirms \cnaclass{ExitingEventArgs} is otherwise unreferenced outside its own
translation unit. This looks like an inconsistency until you check real XNA itself, which
has exactly the same quirk: \texttt{Game.Exiting} is declared \texttt{EventHandler<EventArgs>}
in the real assembly too, despite a dedicated \texttt{ExitingEventArgs} type existing.
CNA preserves this faithfully rather than ``fixing'' it --- a small, concrete example of
Chapter~\ref{ch:design-philosophy}'s ``match FNA, not taste'' rule in action.

\subsection{Properties, and one that is deliberately not settable from outside}

\cnaclass{Game} exposes its state through the project's uniform
\texttt{getXProperty()} / \texttt{setXProperty()} convention: \texttt{Components},
\texttt{Content}, \texttt{GraphicsDevice} (get-only, sourced from the graphics device
service --- there is no \texttt{setGraphicsDeviceProperty}), \texttt{InactiveSleepTime},
\texttt{IsFixedTimeStep}, \texttt{IsMouseVisible}, \texttt{LaunchParameters},
\texttt{TargetElapsedTime}, \texttt{Services}, \texttt{Window}. One property is more
restricted than the rest: \texttt{IsActive} is get-only publicly, with its setter kept
private --- ``only \texttt{Game} itself may change the active state,'' matching FNA's own
\texttt{internal set}. C\texttt{++} has no exact equivalent of C\#'s assembly-internal
visibility, so a private setter reachable only from within \cnaclass{Game}'s own
\texttt{OnActivated} / \texttt{OnDeactivated} handlers is the closest faithful translation.

\subsection{The override points}

A concrete game overrides a fixed set of protected virtuals: \texttt{Initialize()},
\texttt{LoadContent()}, \texttt{UnloadContent()}, \texttt{Update(GameTime\&)},
\texttt{Draw(const GameTime\&)}, \texttt{BeginRun()} / \texttt{EndRun()}, and
\texttt{BeginDraw()} / \texttt{EndDraw()} --- the last of which matters more than it looks:
\texttt{BeginDraw()} returning \texttt{false} skips both \texttt{Draw} and \texttt{EndDraw}
for that frame entirely, the mechanism \cnaclass{GraphicsDeviceManager}
(\S\ref{sec:graphicsdevicemanager-in-loop}) uses to suppress rendering when no device is
available yet. A game can also override \texttt{ShowMissingRequirementMessage(const
std::exception\&)} and the protected \texttt{Dispose(bool disposing)}.

\subsection{The timing loop, ported wholesale from FNA}

\texttt{Game}'s private implementation carries over FNA's adaptive-sleep-precision timing
loop directly: a 128-entry ring buffer of \texttt{previousSleepTimes\_} (masked with
\texttt{SLEEP\_TIME\_MASK} rather than a modulo), a running
\texttt{worstCaseSleepPrecision\_} estimate updated by \texttt{UpdateEstimatedSleepPrecision()},
and \texttt{AdvanceElapsedTime()} driving the actual frame pacing against
\texttt{TargetElapsedTime}, capped by a static \texttt{MaxElapsedTime}. Two pairs of
component-list vectors --- \texttt{updateableComponents\_} / \texttt{currentlyUpdatingComponents\_}
and \texttt{drawableComponents\_} / \texttt{currentlyDrawingComponents\_} --- exist specifically
to guard against a component mutating the component collection (adding or removing itself, or
another component) \emph{during} its own \texttt{Update} / \texttt{Draw} call, a re-entrancy
hazard FNA's own implementation guards against the same way.

One platform-conditional block is CNA-specific rather than ported: under
\texttt{\#if defined(\_\_EMSCRIPTEN\_\_)}, \texttt{Game} gains a private
\texttt{EmscriptenLoopState} struct and an \texttt{EmscriptenMainLoopCallback()}, because a
blocking \texttt{while} loop is not a viable main-loop shape in a browser --- Emscripten
instead requires yielding control back to the browser's own event loop via a registered
callback. This book's Web/Emscripten chapter covers this in full.

\subsection{Two small, explicitly-marked CNA extensions}

Two \texttt{NOXNA}-tagged additions live directly on \cnaclass{Game}: a public
\texttt{RunApplication} flag (documented as ``internal loop flag matching the FNA/XNA
\texttt{Game} implementation shape''), and a pair of compatibility helpers,
\texttt{getTargetFPSProperty()} and \texttt{getTargetMsFrameTimeProperty()}, plus a static
\texttt{fpsToMillisecondsPerFrame(intcs)} --- documented as existing specifically because
``existing CNA examples'' already depend on an FPS-based way to talk about frame timing,
alongside the real XNA \texttt{TargetElapsedTime} \texttt{TimeSpan}.

\section{GameTime: three fields, one privileged writer}

\cnaclass{GameTime} is deliberately small: \texttt{TotalGameTime} and
\texttt{ElapsedGameTime} (both \texttt{TimeSpan}, publicly get-only), and
\texttt{IsRunningSlowly} (bool, publicly get-only). All three setters are \texttt{private},
with \texttt{friend class Game} --- no code outside \cnaclass{Game} itself may mutate a
\cnaclass{GameTime} once constructed, which is exactly the invariant the two-line
\texttt{Update} / \texttt{Draw} signatures from Chapter~\ref{ch:first-game} rely on: a
\cnaclass{GameTime} handed to your \texttt{Update} override is a value your code can read
freely but never falsify.

\section{GameComponent and DrawableGameComponent: the update/draw ordering contract}

\cnaclass{GameComponent} implements \texttt{IGameComponent}, \texttt{IUpdateable}, and
\texttt{System::IComparable<GameComponent>} in one base: a component knows its owning
\cnaclass{Game} (get-only), can be individually \texttt{Enabled}, has an
\texttt{UpdateOrder}, and compares to other components \emph{by} update order via
\texttt{CompareTo} --- which is what lets \cnaclass{GameComponentCollection} keep components
processed in a stable, predictable sequence rather than insertion order.
\cnaclass{DrawableGameComponent} layers \texttt{IDrawable} on top, adding
\texttt{GraphicsDevice} (get-only), \texttt{DrawOrder}, and \texttt{Visible}, and --- notably
--- overrides \texttt{Initialize()} specifically so that content loads lazily once a graphics
device actually exists, via a private \texttt{OnDeviceCreated} handler hooked to graphics
device creation. This mirrors the same ordering constraint from
Chapter~\ref{ch:first-game}: GPU-backed resources can only be created once a device exists,
and \cnaclass{DrawableGameComponent} enforces that automatically for components instead of
leaving it to each component author to get right.

\cnaclass{IGameComponent} itself is minimal --- one pure-virtual \texttt{Initialize()} ---
and \texttt{IUpdateable} / \texttt{IDrawable} are the interface pair that
\cnaclass{GameComponent} / \cnaclass{DrawableGameComponent} implement concretely: get/set
enabled-or-visible state, an ordering value, a changed-event pair, and the actual
\texttt{Update} / \texttt{Draw} method.

\section[GameComponentCollection: an event-raising collection]{GameComponentCollection: an event-raising collection, plus a real C\texttt{++} iterator surface}

\cnaclass{GameComponentCollection} raises \texttt{ComponentAdded} / \texttt{ComponentRemoved}
events (typed \texttt{EventHandler<GameComponentCollectionEventArgs>}) around the usual
collection operations (\texttt{Add}, \texttt{Insert}, \texttt{Remove}, \texttt{RemoveAt},
\texttt{Clear}, \texttt{IndexOf}, \texttt{operator[]}). Because C\# 's
\texttt{IEnumerable<IGameComponent>} has no direct C\texttt{++} equivalent, the collection
adds an explicitly \texttt{NOXNA}-tagged set of \texttt{size\_type} / \texttt{iterator} /
\texttt{const\_iterator} aliases and \texttt{begin()} / \texttt{end()} overloads --- a small,
representative example of how CNA extends an XNA collection type just enough to be usable
with range-based \texttt{for} and the standard library, without touching its XNA-facing
surface.

\section{GameServiceContainer: type-keyed lookup without C\# reflection}

C\#'s \cnaclass{GameServiceContainer} is a dictionary keyed by \texttt{System.Type}. C\texttt{++}
has no runtime reflection to build an equivalent dictionary directly, so CNA's
\texttt{AddService<TService>} / \texttt{GetService<TService>} / \texttt{RemoveService<TService>}
template methods key a \texttt{std::unordered\_map<std::type\_index, void*>} by
\texttt{typeid}, delegating to non-template overloads underneath. Copy construction and copy
assignment are explicitly deleted (``services are registered by pointer identity''), while
move construction and move assignment are defaulted --- a small but precise statement about
ownership: a \cnaclass{GameServiceContainer} can be relocated, but never silently duplicated
into a second set of service pointers.

\subsection{A worked example: registering and consuming a service}

\cnaclass{GameServiceContainer}'s own test suite
(\texttt{GameServiceContainerTests.cpp}) is a good template for the pattern a real game uses
to share a subsystem --- say, a save-game manager --- across otherwise-unrelated components
without threading a pointer through every constructor:

\begin{lstlisting}
struct ISaveGameService {
    virtual void Save(const std::string& slot) = 0;
    virtual ~ISaveGameService() = default;
};

class SaveGameManager final : public ISaveGameService {
public:
    void Save(const std::string& slot) override { /* ... */ }
};

// In Game::Initialize(), or anywhere with access to getServicesProperty():
SaveGameManager saveManager;
getServicesProperty().AddService<ISaveGameService>(&saveManager);

// Anywhere else that only has a GameServiceContainer&:
if (auto* svc = services.GetService<ISaveGameService>()) {
    svc->Save("slot1");
}
\end{lstlisting}

Two behaviors worth knowing before relying on this, both read directly from the container's
own implementation rather than assumed: \texttt{AddService<T>(nullptr)} throws
\texttt{std::invalid\_argument} rather than silently registering an empty slot, and calling
\texttt{AddService<T>} a second time for a type that already has a registered service also
throws \texttt{std::invalid\_argument} --- there is no ``last write wins'' overwrite behavior.
A service must be explicitly \texttt{RemoveService<T>()}'d before a replacement can be
registered under the same type.

\section{GameWindow: one concrete class where FNA has a hierarchy}
\label{sec:gamewindow}

CNA's own class comment for \cnaclass{GameWindow} states the deviation from FNA plainly:
``FNA defines \texttt{GameWindow} as abstract with per-platform subclasses; CNA collapses
that hierarchy into one concrete SDL-backed class.'' Where FNA needs a
\texttt{WinFormsGameWindow}, an SDL2-backed window subclass, and so on, CNA needs exactly one,
because SDL3 already is the cross-platform windowing abstraction --- there is no per-platform
subclassing left to do. \cnaclass{GameWindow} exposes the expected XNA surface
(\texttt{AllowUserResizing}, \texttt{ClientBounds}, \texttt{CurrentOrientation},
\texttt{Handle}, \texttt{ScreenDeviceName}, \texttt{Title}, and the
\texttt{ClientSizeChanged} / \texttt{OrientationChanged} / \texttt{ScreenDeviceNameChanged}
events), plus a small set of explicitly-named \texttt{EXT}-suffixed CNA extensions:
\texttt{GetNativeSdlWindowEXT()} (documented as ``never for use in the strict XNA-facing API
surface itself,'' but needed so e.g.\ \texttt{CNA::Devices::DisplayInfo} can query SDL3
content-scale/safe-area state), \texttt{IsBorderlessEXTProperty}, and
\texttt{MinimizeEXT()} / \texttt{RestoreEXT()} (``CNA extension --- XNA has no minimize/restore
API of its own''). The \texttt{EXT} suffix convention here is a sibling of \texttt{NOXNA}: both
mark non-XNA additions, but \texttt{EXT} specifically marks methods bolted onto an otherwise
XNA-shaped class, kept visually distinct from the XNA members around them.

\section{GraphicsDeviceManager's role in the loop}
\label{sec:graphicsdevicemanager-in-loop}

\cnaclass{GraphicsDeviceManager} --- already seen owning device creation in
Chapter~\ref{ch:first-game} --- implements \cnaclass{IGraphicsDeviceManager}'s two-method
contract (\texttt{BeginDraw() -> bool}, \texttt{CreateDevice()}, plus \texttt{EndDraw()}),
which is exactly what \cnaclass{Game::BeginDraw()} / \texttt{EndDraw()} delegate to: if
\texttt{GraphicsDeviceManager::BeginDraw()} returns \texttt{false} (no device available yet),
\texttt{Game} skips \texttt{Draw} / \texttt{EndDraw} for that frame rather than crashing on a
null device. Its five events (\texttt{DeviceCreated}, \texttt{DeviceDisposing},
\texttt{DeviceReset}, \texttt{DeviceResetting}, \texttt{PreparingDeviceSettings}) and its
protected \texttt{FindBestDevice} / \texttt{RankDevices} / \texttt{CanResetDevice} virtuals
preserve FNA/XNA's device-selection extensibility hooks --- a subclass can override
\texttt{RankDevices} to prefer, say, a higher-resolution adapter, exactly as in real XNA.

One CNA-specific extension lives directly alongside this class rather than inside it: a
\texttt{NOXNA enum class PresentationMode} with five values ---
\texttt{Letterbox}, \texttt{Overscan}, \texttt{Stretch}, \texttt{NativeBackBuffer},
\texttt{FixedHeightDynamicWidth} --- controlling how the game's back buffer is scaled to the
actual display when the two don't match in aspect ratio, exposed through
\texttt{getPreferredPresentationModeProperty()} / \texttt{setPreferredPresentationModeProperty()}.
Real XNA never had to solve this problem in the same way, since the Xbox 360 and a Windows
desktop of that era rarely presented CNA's variety of target displays (a browser canvas, an
Android device in either orientation, an arbitrary desktop window); this is a genuinely new
subsystem, clearly marked as such rather than folded silently into the XNA-facing API.

A game that wants a fixed, letterboxed virtual resolution --- one of the more common real
Getting-Started decisions, since it means every downstream draw call can assume one known
coordinate space regardless of the player's actual window or display size --- sets this before
the device is created, typically in the constructor right alongside \texttt{graphics\_(this)}:

\begin{lstlisting}
MyGame::MyGame()
    : graphics_(this)
{
    graphics_.setPreferredBackBufferWidthProperty(1280);
    graphics_.setPreferredBackBufferHeightProperty(720);
    graphics_.setPreferredPresentationModeProperty(PresentationMode::Letterbox);
}
\end{lstlisting}

With \texttt{Letterbox} selected, a player resizing the window to a wider or taller aspect
ratio than 1280:720 sees black bars rather than a stretched or cropped game --- every draw call
in \texttt{Draw()} keeps working in the same 1280$\times$720 coordinate space CNA's own
internal virtual-resolution machinery (Chapter~\ref{ch:graphicsdevice} covers
\texttt{SetVirtualResolution}, the private method this setting ultimately reaches) scales and
letterboxes on the game's behalf.

\section{Supporting cast: LaunchParameters, TitleContainer, TitleLocation, FrameworkDispatcher}

Four smaller types round out the framework layer. \cnaclass{LaunchParameters} parses process
\texttt{argv} into key/value pairs and --- a CNA-specific representational choice --- inherits
publicly from \texttt{std::unordered\_map<std::string,std::string>} rather than wrapping one,
so a \cnaclass{LaunchParameters} instance \emph{is} a standard associative container as well
as an XNA-shaped one. \cnaclass{TitleContainer} is a static-only class providing
\texttt{OpenStream(path) -> std::unique\_ptr<System::IO::Stream>}, the entry point for
reading title content files portably. \cnaclass{TitleLocation} resolves and caches the base
path content is resolved relative to, exposing it under two accessor names --- both
\texttt{getPathProperty()} and a bare \texttt{Path()} return the same value, the latter kept
specifically ``to match XNA property name'' alongside CNA's own \texttt{getXProperty}
convention. \cnaclass{FrameworkDispatcher}, also static-only, drives framework-level polling
(dynamic audio streams, media, touch input) through a single \texttt{Update()} entry point
that a game is expected to call once per frame; internally it tracks registered
\texttt{Audio::DynamicSoundEffectInstance*} streams behind a mutex --- a thread-safety
addition beyond FNA's own single-threaded assumption.

\section{What this buys a porting engineer}

None of the types in this chapter are exotic on their own. What matters, if you are porting
an existing XNA or FNA game (this book's migration chapter returns to this in practical
detail), is that the entire timing/component/service/window layer a real XNA game's
\texttt{Game}-derived class depends on is present, member-for-member, under the same names.
A \texttt{Game} subclass that only used the public, documented XNA surface --- not FNA's own
internal/assembly-private details --- should port to CNA by translating syntax, not by
redesigning its structure.
